
\section{Templates y Testing}

\subsection{Template y archivos}
\inputcode{cpp}{latex_src/template.h}{template.h}

Ejecutar en la terminal de carpeta
\begin{minted}{bash}
for f in {a..n}; do cp template.cpp $f.cpp; done
\end{minted}

\subsection{Stress testing (diff)}
Para stress testing en problemas donde existe solo una respuesta, usamos $stress.sh$. Testeamos $code.cpp$ contra $brute.cpp$. $gen.cpp$ es el generador de casos. Este ultimo recibe de semilla el valor $i$.

\inputcode{bash}{latex_src/stress.sh}{stress.sh}

Para extraer la semilla en gen.cpp leemos desde argv
\begin{minted}{cpp}
int rnd(int a, int b){return a + rand() % (b - a + 1);}
int main(int argc, char* argv[]){
    int seed = atoi(argv[1]);
    srand(seed);
    int n = rnd(1, 5);
}
\end{minted}


\subsection{Stress testing (checker)}
Usaremos el script similar al anterior $stress.sh$. Aquí el $brute.cpp$ es opcional.

\begin{minted}{cpp}
set -e
g++ code.cpp -o code
g++ gen.cpp -o gen
g++ brute.cpp -o brute
g++ checker.cpp -o checker
for((i = 1; ; ++i)); do
    ./gen $i > input_file
    ./code < input_file > myAnswer
    ./brute < input_file > correctAnswer
    ./checker > checker_log
    echo "Passed test: "  $i
done
\end{minted}

Se ejecutara mientras el return value de checker sea $0$. Ejemplo de la implementación de un checker

\begin{minted}{cpp}
bool validAns(ifstream &fin, int n, int* arr){
    fin>>n;
    //aquí validariamos nuestra respuesta, fin sera el input stream que le pase
    if(algunacosa) return true
    else return false
}
int main(int argc, char * argv[]){
    //input stream del caso
    ifstream fin("input_file", ifstream::in);
    //input stream de mi respuesta
    ifstream ans("myAnswer", ifstream::in);
    //input stream de respuesta bruta
    ifstream cor("correctAnswer", ifstream::in);
    //aqui fin puede procesar lo que interese del input original
    fin>>n;
    int arr[n];
    for(int x = 0; x < n; x++) cin>>arr[x];
    if(!validAns(ans, n, arr)) {
        cout<<"WA x algun motivo, esto saldrá en cheker_log\n";
        return -1;
    }
    return 0;
}
\end{minted}